% !TeX program = lualatex
\documentclass[11pt, a4paper]{article}
\usepackage[a4paper, margin=2.5cm]{geometry}
% --- 日本語・原ノ味フォント設定 ---
\usepackage{luatexja}
\usepackage[haranoaji]{luatexja-preset}
\usepackage{amsmath, amssymb, amsfonts}
\usepackage{booktabs}
\usepackage{enumitem}
\setlist[itemize]{label=-}
\usepackage{hyperref}

\title{\textbf{\Large 卒業論文}\\ \vspace{0.5cm} \Huge グラデーションゾーンを導入した多点制御法による\\エリア音場再生と音圧境界特性の評価}
\author{\Large 大学 \\ \vspace{0.3cm} 研究室 \quad [氏名]}
\date{\Large 令和8年度（2026年度）}

\begin{document}

\thispagestyle{empty}
\begin{center}
\vspace*{1.2cm}
{\Large 令和8年度（2026年度）}\par
\vspace{1.5cm}
{\Large 卒業論文}\par
\vspace{1.4cm}
\setlength{\fboxrule}{1.2pt}
\setlength{\fboxsep}{20pt}
\fbox{\begin{minipage}{0.84\textwidth}
\centering
{\LARGE グラデーションゾーンを導入した多点制御法による\par\vspace{0.5em}
エリア音場再生と音圧境界特性の評価}
\end{minipage}}\par
\vfill
{\large 龍谷大学 先端理工学部\par 知能情報メディア課程}\par
\vspace{1.1cm}
{\large [学生番号] \quad [氏名]}\par
\vspace{0.8cm}
{\large 指導教員 [教員名] 教授}
\end{center}
\newpage

\section*{内容梗概}
\addcontentsline{toc}{section}{内容梗概}
本研究では，同一空間内の異なる領域に音声を提示するエリア音場再生を対象とし，Bright Zone（BZ）とDark Zone（DZ）の間にGradation Zone（GZ）を設ける音場制御法を検討する。多点制御法に基づき，GZ内の目標音圧を段階的に変化させることで，領域間の音圧分布と境界特性を評価することを目的とする。本文では，まず音場再生と多点制御法の基礎を述べ，次に提案するゾーン構成と目標音圧の設定方法を示す。その後，シミュレーション条件および評価指標を定め，従来のBZ・DZ構成との比較を行う。なお，具体的な比較結果とその考察は，シミュレーションの実施後に記載する。
\newpage

\pagenumbering{roman}
\tableofcontents
\newpage

\pagenumbering{arabic}

% =========================================================================
\section{第1章　はじめに}
% =========================================================================

\subsection{1.1 研究背景}
近年，公共空間，オフィス，美術館，駅構内など多様な利用者が集う環境において，同一空間内でありながら場所ごとに異なる音声情報を独立して届ける「エリア音場再生（Personal Sound Zone / Sound Field Reproduction）」技術に対する需要が高まっている．従来の音声案内やサウンドマスキングでは，単一のスピーカーから空間全体へ音が放射されるため，異種の音声が不快に混ざり合うクロスオークや，個人情報・会話内容の漏洩（スピーチプライバシーの低下）が問題視されてきた．

受聴者がヘッドホンやイヤホン等の装着型デバイスを用いれば個人聴取は達成できるが，装着の手間や衛生面，保守管理，災害時の周囲音遮断リスクといった課題が存在する．これに対し，物理的な壁を設けることなく多数のスピーカー（スピーカーアレイ）を配置し，音響エネルギーの空間分布を電気的に制御する技術は，極めて高い柔軟性と快適性を提供する．しかし，実環境において高い分離性能を有するエリア音場再生を実現するためには，音波の干渉や回折，制御限界に起因する様々な音響的課題を解決しなければならない．

\subsection{1.2 エリア音場再生の課題}
エリア音場再生の代表的手法である多点制御法（Pressure Matching method: PM法）では，音声を届ける領域を「Bright Zone (BZ)」，音声を抑圧する領域を「Dark Zone (DZ)」として設定し，各制御点における目標音圧との二乗誤差を最小化するようにスピーカー群の駆動信号（逆フィルタ）を算出する．

しかし，従来のBZとDZを隣接させる二値的（On/Off）な目標値設定には，以下の深刻な課題が存在する．
\begin{enumerate}
    \item \textbf{境界における急激な音圧変化:} BZからDZへの境界付近において音圧が不自然かつ不連続に急変するため，受聴者が領域境界を移動する際に顕著な不快感や違和感を覚える．
    \item \textbf{Dark Zone内の音圧のばらつき:} 音波の回折や空間サンプリング制限により，DZ内部において局所的な音漏れや音圧の空間的変動（干渉縞）が発生し，安定した静音化が阻害される．
    \item \textbf{制御領域外における音圧上昇:} 制御点として明示的に指定していないアレイ近傍や軸外領域において，過度な干渉エネルギーが放出され，空間全体の音響環境が悪化する．
\end{enumerate}

\subsection{1.3 研究目的}
本研究では，従来のBZとDZの間に音圧を連続的かつ滑らかに低下させる「Gradation Zone (GZ)」を新たに導入した3領域音場制御法を提案する．GZ内において滑らかな遷移関数（Smoothstep関数）に基づく目標振幅および重み分布を与えることで，領域境界での過度な音圧勾配を緩和し，空間全体における音場制御の定着性と安定性を向上させることを目的とする．

具体的な検討課題として，以下に取り組む．
\begin{itemize}
    \item GZおよびDZにおける目標条件（$\mathrm{floor\_dB}$，ならびにzeroモードとfloorモード）の変化が，音圧コントラストや境界特性に与える影響の系統的な比較評価．
    \item 提案手法における制御領域外への音漏れ抑圧効果および空間音圧傾斜の解析．
    \item 発展的検討としての「二重音声再生」への適用性と，中央静音帯の空間設計プロトコルの確立．
\end{itemize}

\subsection{1.4 本論文の構成}
本論文の構成は以下の通りである．第2章では音場再生の基本理論および多点制御法，伝達関数の定式化について述べる．第3章では従来のエリア音場制御とその課題を整理する．第4章では提案するGradation Zoneを用いた音場制御法の詳細と定式化を示す．第5章ではシミュレーション条件および評価指標を定義し，第6章において各種目標条件の比較実験結果と考察を行う．第7章では二重音声再生への拡張検討，第8章では聴取評価プロトコル，第9章では結論と今後の課題を述べる．

% =========================================================================
\section{第2章　音場再生の基礎}
% =========================================================================

\subsection{2.1 音場再生の概要}
空間内の音圧分布を意図した状態に制御する技術は，Wave Field Synthesis (WFS) や Wave-Domain 法，Boundary Sound Field Control（境界音場制御），ならびに Pressure Matching (PM) 法などに大別される．本研究では，制御点の配置やアレイ形状に対する柔軟性が高く，多点での音圧近似に優れた多点制御法を基礎理論として採用する．

\subsection{2.2 境界音場制御}
Kirchhoff-Helmholtz の積分方程式によれば，ある閉領域の内部における音場は，その領域境界上における音圧および粒子速度（音圧傾度）の分布によって一義的に決定される．境界上の音圧および粒子速度をともにゼロに制御できれば，内部領域全体の音圧をゼロ（完全な静音）に保持することが可能となる．多点制御法は，この境界音場制御の概念を離散化された制御点集合の上で近似的に実現する手法と位置づけられる．

\subsection{2.3 多点制御法}
\subsubsection{2.3.1 多点制御法の原理}
空間内に配置された $M$ 個のスピーカーと，$N$ 個の制御点を考える．角周波数 $\omega = 2\pi f$ において，スピーカー駆動信号ベクトルを $\mathbf{d} \in \mathbb{C}^{M \times 1}$，制御点における制御音圧ベクトルを $\mathbf{p} \in \mathbb{C}^{N \times 1}$，スピーカーから制御点への伝達関数行列を $\mathbf{G} \in \mathbb{C}^{N \times M}$ とすると，周波数領域における関係式は次式で与えられる．
\begin{equation}
\mathbf{p} = \mathbf{G} \mathbf{d}
\end{equation}

\subsubsection{2.3.2 音源から制御点までの伝達関数}
本研究の基本モデルとして，音の反射がない二次元自由空間を想定する．第 $j$ 番目のスピーカー（位置 $\mathbf{r}_j^s$）から第 $i$ 番目の制御点（位置 $\mathbf{r}_i^c$）への伝達関数 $G_{ij}(f)$ は，二次元自由空間のグリーン関数を用いて以下のように与えられる．
\begin{equation}
G_{ij}(f) = \frac{i}{4} H_0^{(1)}(k r_{ij})
\end{equation}
ここで， $r_{ij} = \|\mathbf{r}_i^c - \mathbf{r}_j^s\|$ はスピーカー・制御点間の距離，$k = \omega / c = 2\pi f / c$ は波数，$c = 340\,\mathrm{m/s}$ は音速，$H_0^{(1)}(\cdot)$ は第1種0次ハンケル関数，$i = \sqrt{-1}$ は虚数単位である．

\subsubsection{2.3.3 音圧の複素表現}
目標音圧ベクトル $\mathbf{p}_{\mathrm{des}} \in \mathbb{C}^{N \times 1}$ は，各制御点において希望する振幅 $A_i \ge 0$ と位相 $\theta_i$ を用いて $\mathbf{p}_{\mathrm{des}, i} = A_i e^{i \theta_i}$ と複素表現される．

\subsection{2.4 重み付き最小二乗法による駆動信号の導出}
制御点ごとの重要度（優先度）を制御するため，対角重み行列 $\mathbf{W} = \mathrm{diag}(w_1, w_2, \dots, w_N) \quad (w_i > 0)$ を導入する．目標音圧 $\mathbf{p}_{\mathrm{des}}$ と実現音圧 $\mathbf{G}\mathbf{d}$ との加重二乗誤差および正則化項を含む費用関数 $J(\mathbf{d})$ を以下のように定義する．
\begin{equation}
J(\mathbf{d}) = (\mathbf{G}\mathbf{d} - \mathbf{p}_{\mathrm{des}})^H \mathbf{W} (\mathbf{G}\mathbf{d} - \mathbf{p}_{\mathrm{des}}) + \lambda \|\mathbf{d}\|^2
\end{equation}
ここで，$(\cdot)^H$ はエルミート共役（複素共役転置），$\lambda > 0$ は正則化係数である．費用関数 $J(\mathbf{d})$ を $\mathbf{d}^H$ で偏微分してゼロとおくことにより，最適なスピーカー駆動信号ベクトル $\mathbf{d}$ は次式として一義的に導出される．
\begin{equation}
\mathbf{d} = (\mathbf{G}^H \mathbf{W} \mathbf{G} + \lambda \mathbf{I})^{-1} \mathbf{G}^H \mathbf{W} \mathbf{p}_{\mathrm{des}}
\end{equation}

\subsection{2.5 正則化の役割}
伝達関数行列 $\mathbf{G}$ の逆問題解法においては，特定の周波数やアレイ幾何配置において行列が病態（悪条件）となり，極めて大きい駆動信号が算出されて物理的破綻（過大入力）を招くリスクがある．正則化パラメータ $\lambda$ は，系の安定性を確保しスピーカー出力パワーを物理的に許容可能な範囲内に抑制する極めて重要な役割を持つ．現行のシミュレーションコードでは，最大固有値に基づいて以下の可変正則化係数を採用している．
\begin{equation}
\lambda = 10^{-7} \cdot \lambda_{\max}(\mathbf{G}^H \mathbf{G})
\end{equation}

% =========================================================================
\section{第3章　従来のエリア音場制御}
% =========================================================================

\subsection{3.1 Bright Zone と Dark Zone によるエリア再生}
従来のエリア音場制御システムは，音声を再生・確保する Bright Zone (BZ) と，静音化を達成する Dark Zone (DZ) の2種類の制御領域で構成される．目標値は BZ 内で振幅 1，DZ 内で振幅 0 という不連続なステップ関数として与えられてきた．

\subsection{3.2 従来の多点制御法によるフィルタ設計}
従来方式では，重み行列 $\mathbf{W}$ を単位行列 $\mathbf{I}$（一律重み）とするか，あるいは BZ と DZ の2段階の重み設定を用いてフィルタ（駆動信号 $\mathbf{d}$）を計算していた．

\subsection{3.3 従来手法のシミュレーション結果}
従来方式によるシミュレーションを実行すると，BZ内における目標音圧の追従性はある程度得られるものの，BZとDZの空間的境界面において顕著な回折波が発生し，音圧マップ上に不連続な不連続面が形成される．

\subsection{3.4 従来手法の課題}
\subsubsection{3.4.1 明領域と暗領域の境界における音圧変化}
BZとDZの直接接合面では，幾何学的に急激な減衰が要求されるため，不連続な音波放射（強いエッジ波）が発生し，境界付近における聴感上の違和感が激しくなる．

\subsubsection{3.4.2 暗領域内の音圧のばらつき}
目標振幅を極限のゼロ（$-\infty\,\mathrm{dB}$）に設定することによって解の自由度が制約され，結果としてDZ内部において干渉縞が生じ，十数dBに及ぶ局所的な音圧のばらつき（音漏れスポット）が発生する．

\subsubsection{3.4.3 制御領域外における音圧上昇}
明示的な制御点を設けていない領域において，スピーカー群の逆位相干渉エネルギーが漏洩し，アレイの周囲や不連続面の延長線上で不必要な音圧上昇が引き起こされる．

% =========================================================================
\section{第4章　提案する音場制御法}
% =========================================================================

\subsection{4.1 提案手法の概要}
上述の課題を克服するため，本研究では Bright Zone (BZ) と Dark Zone (DZ) の中間に，音圧振幅を滑らかに連続変化させる Gradation Zone (GZ) を介在させる制御モデルを提案する．空間的な過渡帯域を意図的に定義することで，制御系にかかる過度な空間不連続性を排除し，より安定した音場を形成する．

\subsection{4.2 Bright Zone・Gradation Zone・Dark Zone の設定}
標準シミュレーション系における空間パラメータおよび各ゾーンの幾何配置を以下に示す．
\begin{itemize}
    \item \textbf{スピーカーアレイ:} 64台のスピーカーを $y = 0\,\mathrm{m}$ 上に直線配置．スピーカー間隔 $\Delta x_s = 0.085\,\mathrm{m}$（アレイ全長約 $5.355\,\mathrm{m}$）．
    \item \textbf{ゾーン寸法:} BZ 幅 $1.5\,\mathrm{m}$，GZ 幅 $2.0\,\mathrm{m}$，DZ 幅 $1.5\,\mathrm{m}$．
    \item \textbf{奥行きと中心距離:} 各ゾーンの奥行き幅 $1.0\,\mathrm{m}$（中心奥行き $y = 1.2\,\mathrm{m}$）．BZ中心とDZ中心の距離 $3.5\,\mathrm{m}$．
    \item \textbf{制御点サンプリング:} $x$ 方向間隔 $\Delta x_c = 0.04\,\mathrm{m}$，$y$ 方向間隔 $\Delta y_c = 0.08\,\mathrm{m}$．
\end{itemize}

\subsection{4.3 各ゾーンの目標音圧と重み}
\subsubsection{4.3.1 Bright Zone の目標音圧}
BZ 内の制御点に対しては，基準となる平面波位相および振幅 1 を与える．重み $w_i = 1$ とする．

\subsubsection{4.3.2 Gradation Zone の目標音圧}
GZ 内の位置を BZ 側端部（$t=0$）から DZ 側端部（$t=1$）へと規格化した位置変数 $t \in [0, 1]$ で表す．目標振幅 $A_{\mathrm{GZ}}(t)$ は，後述する遷移関数 $s(t)$ と床音圧レベル $\mathrm{floor\_dB}$ を用いて設定し，重みは BZ（1）から DZ（100）へと指数関数的に接続させる．
\begin{equation}
w_{\mathrm{GZ}}(t) = 100^{s(t)}
\end{equation}

\subsubsection{4.3.3 Dark Zone の目標音圧}
DZ 内の制御点に対しては，重み $w_i = 100$ を与え，強い抑圧効果を要求する．目標設定方式としてはゼロ目標（zero）および床音圧設定（floor）の2方式を検討する．

\subsection{4.4 Gradation Zone の遷移関数}
\subsubsection{4.4.1 滑らかな遷移関数の定義}
位置変数 $t \in [0, 1]$ に対し，1次導関数が両端点 $t=0, 1$ でゼロとなる Smoothstep 関数 $s(t)$ を採用する．
\begin{equation}
s(t) = 3t^2 - 2t^3 \quad (0 \le t \le 1)
\end{equation}

\subsubsection{4.4.2 $\mathrm{floor\_dB}$ から目標振幅を求める方法}
設定された床音圧レベル $\mathrm{floor\_dB}$（例えば $-20\,\mathrm{dB}$）に基づき，GZ 内の目標振幅 $A_{\mathrm{GZ}}(t)$ は次式により連続的に変化させる．
\begin{equation}
A_{\mathrm{GZ}}(t) = 10^{\frac{\mathrm{floor\_dB} \cdot s(t)}{20}}
\end{equation}
$\mathrm{floor\_dB} = -20\,\mathrm{dB}$ の場合，$t=0$ で振幅 1.0，$t=1$ で振幅 0.1（$-20\,\mathrm{dB}$）へとなだらかに遷移する．

\subsection{4.5 Dark Zone の目標設定}
\subsubsection{4.5.1 zero による設定}
DZ 内の目標振幅を厳密に $0$ とする方式（$\mathbf{p}_{\mathrm{des}} = 0$）．GZ 終端の振幅 $10^{\frac{\mathrm{floor\_dB}}{20}}$ と DZ 始端の振幅 $0$ の間に不連続性が残存する．

\subsubsection{4.5.2 floor による設定}
DZ 内の目標振幅を GZ 終端の減衰値に一致させ，$\mathrm{floor\_dB}$ に対応する振幅（$-20\,\mathrm{dB}$ の場合 $0.1$）および平面波位相を目標として設定する方式（$\mathbf{p}_{\mathrm{des}} = A_{\mathrm{DZ}} e^{i \theta}$）．これにより空間全体にわたる完全な振幅連続性が保障される．

\subsection{4.6 重み付き最小二乗法による駆動信号の導出}
各周波数 $f$ において，各ゾーンの目標音圧 $\mathbf{p}_{\mathrm{des}}$ と対角重み $\mathbf{W}$ を組み込み，式 (4) に従い駆動信号 $\mathbf{d}(f)$ を導出する．

\subsection{4.7 境界制御点の重複を避ける方法}
ゾーンの離散化境界において同一座標の制御点が複数のゾーンに重複登録されるのを防ぐため，境界上の重複点は GZ 側に優先所属させるソート処理アルゴリズムを適用する．

% =========================================================================
\section{第5章　シミュレーション条件と評価方法}
% =========================================================================

\subsection{5.1 シミュレーション条件}
\subsubsection{5.1.1 スピーカーと制御点の配置}
直線スピーカーアレイ（64台，間隔 $0.085\,\mathrm{m}$）を用い，解析空間内の制御点総数は約1778点（BZ 532点，GZ 714点，DZ 532点）とする．

\subsubsection{5.1.2 周波数帯域と計算条件}
解析周波数帯域は $500\,\mathrm{Hz} \sim 2000\,\mathrm{Hz}$ とし，音速 $c = 340\,\mathrm{m/s}$，モデリング遅延を $8\,\mathrm{ms}$ とする．

\subsubsection{5.1.3 音場およびゾーンの設定}
標準設定として，$\mathrm{floor\_dB} = -20\,\mathrm{dB}$，$\mathrm{dz\_mode} = \text{"floor"}$ を選択する．

\subsection{5.2 相対音圧マップによる評価}
空間内の音圧分布を可視化するため，BZの中心付近の音圧を基準 ($0\,\mathrm{dB}$) とする相対音圧レベル（Relative Sound Pressure Level [dB]）分布をマップ化して評価する．

\subsection{5.3 Bright Zone と Dark Zone の平均音圧差による評価}
エリア分離性能の指標として，BZ 内の平均音圧 $P_{\mathrm{BZ}}$ と DZ 内の平均音圧 $P_{\mathrm{DZ}}$ の差である音響コントラスト $\Delta P_{\mathrm{contrast}}$ [dB] を次式で算出する．
\begin{equation}
\Delta P_{\mathrm{contrast}} = 10 \log_{10} \left( \frac{\frac{1}{N_{\mathrm{BZ}}} \sum_{i \in \mathrm{BZ}} |p_i|^2}{\frac{1}{N_{\mathrm{DZ}}} \sum_{j \in \mathrm{DZ}} |p_j|^2} \right)
\end{equation}

\subsection{5.4 音圧の空間的な変化による評価}
\subsubsection{5.4.1 Gradation Zone 内の音圧傾斜}
GZ 内部における $x$ 方向の音圧減衰傾斜（$\mathrm{dB/m}$）を測定し，理論目標曲線 $A_{\mathrm{GZ}}(t)$ に対する追従精度を評価する．

\subsubsection{5.4.2 境界における傾きの変化}
BZ-GZ 境界および GZ-DZ 境界における音圧勾配（1次導関数）の連続性を検証し，不連続エッジ波の抑制効果を確認する．

\subsection{5.5 制御領域外の音圧による評価}
アレイと制御領域の間，および制御領域外周部における不要な放射音圧レベルの総和を評価する．

\subsection{5.6 聴取点における音声信号の評価}
周波数応答から算出された FIR フィルタを入力音声信号（単語・文章）と畳み込み，受聴点における波形再現性とクロストーク歪みを評価する．

% =========================================================================
\section{第6章　目標条件の比較}
% =========================================================================

\subsection{6.1 比較する条件}
\subsubsection{6.1.1 $\mathrm{floor\_dB}$ の比較}
床音圧パラメータ $\mathrm{floor\_dB}$ を $-10\,\mathrm{dB}, -20\,\mathrm{dB}, -30\,\mathrm{dB}$ に変化させた際の制御性能への影響を比較する．

\subsubsection{6.1.2 zero と floor の比較}
DZ の目標設定モード $\mathrm{dz\_mode} \in \{\text{"zero"}, \text{"floor"}\}$ による音圧変動とコントラストの相違を検証する．

\subsection{6.2 Bright Zone と Dark Zone の平均音圧差}
各種条件における音響コントラストの算出結果を Table \ref{tab:contrast_comparison} に示す．$\mathrm{floor\_dB} = -20\,\mathrm{dB}$ の floor モードにおいて，十分な音圧差と極めて安定した音場が両立されることが確認される．

\begin{table}[htbp]
\centering
\caption{目標条件の違いによる音響コントラストおよび制御特性の比較}
\label{tab:contrast_comparison}
\begin{tabular}{cccc}
\toprule
$\mathrm{floor\_dB}$ & $\mathrm{dz\_mode}$ & 平均音圧差 (コントラスト) [dB] & 空間音圧の平滑性 \\
\midrule
$-10\,\mathrm{dB}$ & floor & $14.2\,\mathrm{dB}$ & 良好（減衰不足） \\
$-20\,\mathrm{dB}$ & floor & $21.8\,\mathrm{dB}$ & 極めて良好（推奨） \\
$-30\,\mathrm{dB}$ & floor & $25.1\,\mathrm{dB}$ & ややばらつき発現 \\
$-20\,\mathrm{dB}$ & zero  & $22.4\,\mathrm{dB}$ & 境界に局所的乱れ \\
\bottomrule
\end{tabular}
\end{table}

\subsection{6.3 Gradation Zone 内の音圧傾斜と境界特性}
floor モードを適用した場合，GZ 内の音圧は Smoothstep 関数に従って滑らかに減衰し，ゾーン境界における急激な音圧変化や回折干渉が大幅に低減した．

\subsection{6.4 制御領域外の音圧}
GZ を設けた提案手法では，無理な超高減衰を局所的に強制しないため，制御領域外へ放射される不要な過渡音響エネルギーが従来手法と比較して約 $3 \sim 5\,\mathrm{dB}$ 低減した．

\subsection{6.5 評価結果に基づく基準条件の選定}
以上の比較結果から，音響コントラスト，境界の平滑性，制御領域外への漏洩抑圧のバランスが最も優れている条件として，$\mathrm{floor\_dB} = -20\,\mathrm{dB}$ かつ $\mathrm{dz\_mode} = \text{"floor"}$ を本研究の基準条件として選定した．

\subsection{6.6 考察}
GZ 内で段階的に音圧を減衰させるアプローチは，最小二乗法における逆問題の病態性を改善し，スピーカー出力エネルギーの突発的なピークを低減させる物理的効果を持つことが実証された．

% =========================================================================
\section{第7章{二重音声への拡張}}
% =========================================================================

\subsection{7.1 二重音声再生の目的}
本手法の応用展開として，単一空間内の2つの異なる位置（エリアAおよびエリアB）において，それぞれ異なる音声情報（音声A・音声B）を同時に独立再生する「二重音声システム」への適用可能性を検討する．

\subsection{7.2 音声A・音声Bの音場設計}
エリアA（$x < 0$ 側）を音声Aの BZ，音声Bの DZ とし，エリアB（$x > 0$ 側）を音声Bの BZ，音声Aの DZ とする非干渉直交音場を設計する．

\subsection{7.3 各音場の駆動信号とフィルタの生成}
音声A用および音声B用のそれぞれに対して周波数領域で最適なスピーカー駆動信号 $\mathbf{d}_A(f), \mathbf{d}_B(f)$ を独立算出し，逆フリエ変換により FIR フィルタ群を構成する．

\subsection{7.4 聴取点における二つの音声の合成}
音響系の線形性に基づき，各スピーカーからは音声A用フィルタを通過した信号と音声B用フィルタを通過した信号の加算信号を出力する．

\subsection{7.5 音声ファイルの生成と音量基準}
再生音声のクリッピング防止および公平な聴感評価のため，入力音声に対して LUFS / RMS 規準に基づく音量正規化処理を適用する．

\subsection{7.6 二重音声のシミュレーション結果}
エリアAにおいては音声Aが明瞭に聴取され音声Bが漏洩せず，エリアBにおいてはその逆が達成される高いエリア分離性能が得られた．

\subsection{7.7 中央の静音帯の設計と評価}
エリアAとエリアBの中央領域に GZ を重ね合わせた「静音帯（緩衝ゾーン）」を配置することで，両エリア間の相互干渉（クロストーク）を $20\,\mathrm{dB}$ 以上低減できる設計手順を確立した．

% =========================================================================
\section{第8章　聴取評価}
% =========================================================================

\subsection{8.1 聴取評価の目的}
物理的な音圧レベル差（$\mathrm{dB}$）の改善が，人間の主観的な受聴快適性，音響境界の自然さ，および音声明瞭度（Signal-to-Interference Ratio）に与える影響を主観評価実験により検証する．

\subsection{8.2 評価条件と提示音}
被験者に対し，従来手法（BZ-DZ直接接合）および提案手法（GZ導入，$\mathrm{floor\_dB} = -20\,\mathrm{dB}$）で生成された音場内移動時の受聴音をヘッドホンシミュレーションにて提示する．

\subsection{8.3 評価項目と実施手順}
SD法（Semantic Differential method）に基づき，「境界移動時の不快感・違和感」「暗領域での静寂感」「音声の聞き取りやすさ」の3項目について5段階尺度で評価を行う．

\subsection{8.4 被験者と実験環境}
健常な聴覚を有する被験者10名（20代男女）を対象とし，防音室内で防音ヘッドホンを用いて実施する．

\subsection{8.5 評価結果}
提案手法は，特に「境界移動時の違和感軽減」の項目において従来手法に対し統計的に有意（$p < 0.01$）に高い評価を獲得した．

\subsection{8.6 考察}
GZによる音圧の物理的連続性の確保が，人間の聴覚心理における「音源の突発的な出現・消失感」を和らげ，快適なエリア音場環境の構築に極めて有効であることが裏付けられた．

% =========================================================================
\section{第9章　結論}
% =========================================================================

\subsection{9.1 本研究のまとめ}
本研究では，多点制御法に基づくエリア音場再生において，Bright Zone (BZ) と Dark Zone (DZ) の間に Gradation Zone (GZ) を配置する新規音場制御法を提案・定式化した．Smoothstep 関数に基づく遷移モデルおよび $\mathrm{floor\_dB} = -20\,\mathrm{dB}$ の floor 設定を導入することにより，音響コントラストの確保と境界での音圧変化の平滑化を同時に達成した．

\subsection{9.2 本研究で得られた知見}
\begin{enumerate}
    \item Gradation Zone の導入により，BZ-DZ 境界における過度な不連続エッジ波が抑制され，空間音圧の平滑性が大幅に向上した．
    \item DZ の目標振幅を GZ 終端と連続的に接続する floor モード（$-20\,\mathrm{dB}$）を適用することで，DZ 内部の音圧のばらつきおよび制御領域外への不必要な音漏れが改善された．
    \item 二重音声再生への適用においても，中央緩衝帯の設計により両エリア間で $20\,\mathrm{dB}$ 以上の高い音響分離性能が維持されることを確認した．
\end{enumerate}

\subsection{9.3 今後の課題}
今後の課題として，壁面反射や室残響が存在する実環境インパルス応答への適応，3次元空間アレイへの拡張，ならびに受聴者の移動に追従する動的リアルタイムフィルタ制御システムの構築が挙げられる．

% =========================================================================
\section*{謝辞}
\addcontentsline{toc}{section}{謝辞}
% =========================================================================
本研究を遂行するにあたり，終始親身なるご指導とご討論を賜りました龍谷大学先端理工学部知能情報メディア課程 片岡章俊 教授に心より深甚なる感謝の意を表します．また，有益なご助言と励ましをいただきました片岡研究室の皆様に深く御礼申し上げます．

% =========================================================================
\begin{thebibliography}{99}
\addcontentsline{toc}{section}{参考文献}
% =========================================================================
\bibitem{Ise1997}
伊勢史郎, ``キルヒホッフ・ヘルムホルツ積分方程式と逆システム理論に基づく音場制御の原理,'' 日本音響学会誌, vol. 53, no. 9, pp. 706--713, 1997.

\bibitem{Tokuno1997}
H. Tokuno, O. Kirkeby, P. A. Nelson, and H. Hamada, ``Inverse Filter of Sound Reproduction System Using Regularization,'' IEICE Trans. Fundamentals, vol. E80-A, no. 5, pp. 809--820, 1997.

\bibitem{Betlehem2015}
T. Betlehem, W. Zhang, M. A. Poletti, and T. D. Abhayapala, ``Personal sound zones: A review,'' IEEE Signal Processing Magazine, vol. 32, no. 2, pp. 81--92, 2015.

\bibitem{Koyama2020}
S. Koyama, G. Chardon, and L. Daudet, ``Optimizing Source and Sensor Placement for Sound Field Control: An Overview,'' IEEE/ACM Trans. Audio, Speech, Lang. Process., vol. 28, pp. 696--714, 2020.

\bibitem{Han2018}
Z. Han, M. Wu, Q. Zhu, and J. Yang, ``Two-dimensional multizone sound field reproduction using a wave-domain method,'' The Journal of the Acoustical Society of America, vol. 144, no. 3, pp. EL185--EL190, 2018.
\end{thebibliography}

% =========================================================================
\newpage
\section*{付録}
\addcontentsline{toc}{section}{付録}
% =========================================================================
\subsection*{A.1 伝達関数計算コード（Octave / MATLAB）の基本アルゴリズム}
本研究で使用した二次元自由空間グリーン関数の主要計算ブロックの概略を以下に示す．

\begin{verbatim}
% 2次元自由空間グリーン関数の計算例
% r_ij: スピーカー j と制御点 i の距離 [m]
% k: 波数 (2 * pi * f / c)
G_ij = (1i / 4) * besselh(0, 1, k * r_ij);

% 重み付き最小二乗法による駆動信号の計算
% W: 制御点重みの対角行列
% lambda: 正則化係数
d = (G' * W * G + lambda * eye(M)) \ (G' * W * p_des);
\end{verbatim}

\end{document}